\section{Bounded elementary number theory}
\label{sec:elementary}

\begin{definition}[Primality]
For a natural number $p$, $\Prime(p)$ means
\[
 p\ge2\quad\text{and}\quad
 \forall d<p\;\bigl(2\le d\to p\bmod d\ne0\bigr).
\]
This notation abbreviates a primitive recursive Boolean test. Divisibility
by a positive natural is also a remainder-zero test and supplies the
quotient as its witness.
\end{definition}

\begin{lemma}[A bounded least divisor is prime]
\label{lem:lpf}
For $q\ge2$, scan the finite list $2,3,\ldots,q$ and let $p$ be its first
divisor of $q$. This search succeeds, $2\le p\le q$, $p\mid q$, and $p$ is
prime. All the functions and correctness assertions are available in $\ipra$.
\end{lemma}

\begin{proof}
The last candidate $q$ divides $q$, so the scan cannot fail. If
$2\le d<p$ and $d\mid p$, then $d\mid q$: multiply the exact quotient
identities for $p/d$ and $q/p$. That contradicts the recorded failure of
every candidate smaller than $p$. Each candidate is a decidable test, and
the scan's correctness is proved by quantifier-free induction on its index.
No assertion about existence or uniqueness of a full prime factorization
is used.
\end{proof}

\begin{lemma}[Prime divisibility of a product]
\label{lem:euclid}
If $p$ is prime and $p\mid uv$ for signed integers $u,v$, then
$p\mid u$ or $p\mid v$. This disjunction has a constructive $\ipra$ proof.
\end{lemma}

\begin{proof}
Compute $a=u\bmod p$, so $0\le a<p$. If $a=0$, the first disjunct holds.
Assume $a>0$. We give the bounded arithmetic proof that $a$ has a
B\'ezout relation with $p$.

Initialize remainders and signed coefficients by
\[
 r_0=p=1p+0a,\qquad r_1=a=0p+1a.
\]
At an active stage, with $r_1>0$, compute
$Q=r_0\mathbin{\mathrm{div}}r_1$ and $r_2=r_0-Qr_1$ and replace
$(r_0,r_1)$ by $(r_1,r_2)$. Replace their coefficient pairs in exactly the
same way. If $r_1=0$, retain the entire record. Run exactly $p+1$ stages.

At every active stage $0\le r_2<r_1$; hence the second remainder decreases
by at least one until it is zero. Quantifier-free induction gives the bound
$r_1(t)\le a\monus t$. After $p+1$ stages it is zero. The first remainder
stays positive, and both remainders stay at most $p$. The coefficient
identities are preserved because the same linear operation is applied to
both sides of each identity.

For every $1\le d\le p$, the relation
\[
 (d\mid r_0\land d\mid r_1)
 \quad\Longleftrightarrow\quad
 (d\mid r_1\land d\mid r_0-Qr_1)
\]
follows by substitution of the quotient identities, in each direction.
Consequently the preserved common-divisor invariant is the finite conjunction
\[
 \bigwedge_{1\le d\le p}
 \bigl((d\mid r_0(t)\land d\mid r_1(t))
       \leftrightarrow(d\mid p\land d\mid a)\bigr).
\]
Represent this conjunction by its Boolean fold. It holds initially, and
the displayed one-step equivalence proves its preservation by transitivity,
using quantifier-free induction. At the final record $(g,0)$,
it follows, by taking $d=g$, that $g\mid p$ and $g\mid a$.
Here $1\le g\le p$.

Primality gives $g=1$ or $g=p$: a value $2\le g<p$ is excluded directly
by the primality test. The latter alternative is impossible because
$0<a<p$. Thus $g=1$. The computed coefficient record supplies actual
integers $s,t$ with
\[
 sp+ta=1.
\]
The coefficient update is primitive recursive. For an optional explicit
size bound, every quotient is at most $p$, and the maximum absolute
coefficient grows by a factor of at most $p+1$ per stage. Thus
$(p+1)^{p+1}$ bounds its size after $p+1$ stages.

Because $u\equiv a\pmod p$, the hypothesis gives $p\mid av$.
Multiplying $sp+ta=1$ by $v$ shows that $p\mid v$. The initial test
$a=0$ decides which disjunct is obtained. Every loop has a prescribed
bound and a coded quantifier-free invariant, so the proof uses only
allowed induction.
\end{proof}

\section{A finite modular seed for every prime}
\label{sec:seed}

The prime $2$ already has the tuple $(1,1,0,0)$. A prime greater than $2$
is odd, since an even number greater than $2$ has the forbidden divisor
$2$. Hence write an odd prime as $p=2h+1$, with $h\ge1$.

\begin{lemma}[Two short lists of square residues]
\label{lem:residue}
For $p=2h+1$ prime, there are explicitly computable $a,b\le h$ with
\[
 a^2+b^2+1\equiv0\pmod p.
\]
One may find them by at most $(h+1)^2$ tests.
\end{lemma}

\begin{proof}
For $0\le i,j\le h$, form the residues in $\{0,\ldots,p-1\}$
\[
 A_i=i^2\bmod p,\qquad B_j=(-1-j^2)\bmod p.
\]
Each list is injective. If $A_i=A_j$, then
$p\mid(i-j)(i+j)$. Lemma~\ref{lem:euclid} gives $p\mid i-j$ or
$p\mid i+j$. The bounds $\abs{i-j}<p$ and $0\le i+j\le2h=p-1$
then give $i=j$ in the first case, or $i=j=0$ in the second.
Indeed an integer divisible by $p$ and of absolute value less than $p$
must be zero, by the exact quotient equation. Equality of two $B$ entries
reduces to the same square-residue argument.

Scan the $(h+1)^2$ pairs in lexicographic order, storing the first pair
with $A_a=B_b$ and a hit flag. We prove that the failure flag is impossible
by an explicit finite counting argument. If there were no cross-list
match, the concatenation $L$ of the two lists would have
$2(h+1)=p+1$ pairwise distinct entries, each smaller than $p$.

For $r<p$, let
\[
 C(r)=\sum_{k=0}^{p}\ind{L_k=r},
\]
where the bracket is the equality-test bit. Because the entries are
pairwise distinct, at most one summand is $1$, so $C(r)\le1$.
For each $k$, exactly one of the $p$ equality tests $\ind{L_k=r}$ is $1$,
because $L_k<p$. Counting the finite incidences in the two orders gives
\begin{equation}
\label{eq:count}
 p+1
 =\sum_{k=0}^{p}\sum_{r=0}^{p-1}\ind{L_k=r}
 =\sum_{r=0}^{p-1}C(r)
 \le p,
\end{equation}
a contradiction. The interchange of the two finite sums is proved by
successively adding rows and distributing addition. These are nested
quantifier-free inductions on the two finite bounds. The assertion that
at most one summand is $1$ can either be recorded as a Boolean prefix
condition or proved by induction on the count; no set cardinalities or
infinite combinatorial principles are involved.

The hit flag is a computed bit. Its $0$ case has just been excluded, so it
is $1$ and its stored pair satisfies $A_a=B_b$. This equality is precisely
the desired congruence. The scan's index $k<(h+1)^2$ can be decoded as
\[
 a=k\mathbin{\mathrm{div}}(h+1),\qquad b=k\bmod(h+1),
\]
which specifies the search as a single primitive recursive bounded loop.
\end{proof}

\begin{proposition}[Small positive multiple of a prime]
\label{prop:seed}
For each prime $p$ a primitive recursive seed construction returns
$m_0\in\NN$ and $x^{(0)}\in\ZZ^4$ with
\[
 1\le m_0<p,\qquad \Nrm{x^{(0)}}=m_0p.
\]
For odd $p$ it can ensure $m_0\le(p-1)/2$.
\end{proposition}

\begin{proof}
For $p=2$ take $m_0=1$ and $x^{(0)}=(1,1,0,0)$.
For odd $p=2h+1$, use the successful pair of Lemma~\ref{lem:residue} and set
\[
 x^{(0)}=(a,b,1,0),\qquad
 m_0=(a^2+b^2+1)\mathbin{\mathrm{div}}p.
\]
The division is exact. Its numerator is positive, so $m_0\ge1$.
Since $h\ge1$,
\[
 a^2+b^2+1\le2h^2+1\le2h^2+h=hp.
\]
Cancellation of the positive factor $p$ gives $m_0\le h<p$.
For nonprime inputs the same finite program is made total by returning
an arbitrary default, namely $(1,(0,0,0,0))$, when no pair is found.
The theorem asserts correctness only under the decidable prime condition.
\end{proof}
